
\section{Background}


\textbf{LLM Scaling Laws.} LLMs are typically decoder-only transformers \citep{vaswani2017attention} trained via next token prediction on web-scale text \citep{radford2019language}. It has been empirically observed that the performance of LLMs scale well with model size -- with the largest model showing emergent capabilities \citep{wei2022emergent}. \cite{kaplan2020scalinglaw} shows that LLM performance roughly scales as a \textbf{power-law} in the model size $N$. The performance here is measured by validation loss $L$, which is well known to correlate strongly with downstream metrics. Specifically they propose the following law (using constants $N_c, \alpha_N$) for models trained on sufficiently large datasets: $L(N) = (N_c/N)^{\alpha_N}$. \citet{hoffmann2022training} also showed that the performance scales well with \textit{token horizon} and that the optimal performance for a given FLOPs budget is obtained by scaling the model and token horizon jointly. The current paradigm thus scales token count in addition to model size -- e.g. LLama-3 was trained on $\textgreater$10x as many tokens as LLama-1 \citep{llama3}. In this paper we use notation from \cite{kaplan2020scalinglaw}, denoting the token horizon (i.e. number of tokens seen during training, some of which might be duplicate) by $D$ and the model size (i.e. parameter count) by $N$. As is common in the literature, we will fit scaling laws to empirical observations. \rebuttal{Following \cite{kaplan2020scalinglaw}, we will use power-laws of the form $F(X) = A X^\alpha$, where $\alpha$ and $A$ are fitted and $X$ is the quantity of interest.} To measure goodness-of-fit we will use the $R^2$ measure -- its value ranges from 0 to 1, where 1 indicates a perfect fit and 0 indicates no fit. 





\textbf{Hyperparameter transfer.} Hyperparameters can strongly influence the performance of LLMs, and LR is a critical hyperparameter. The muP paper \citep{mup} first popularized \textit{hyperparameter transfer} -- the process of finding optimal hyperparameters from a small proxy experiment. Core to muP is the muP-parametrization -- a way of parametrizing an LLM such that the optimal LR for a small model is also optimal for a larger model. The muP parametrization introduces some changes to the network, e.g. the attention scaling factor is changed from $\frac{1}{\sqrt{d}}$ to $\frac{1}{d}$. The original muP paper shows that LR and many other hyperparameters transfer from small to large models when muP parametrization is used, which is further corroborated in \cite{testingmup}. While \cite{mup} focuses on transfer between model sizes, we focus on transfer between token horizons.



